\documentclass{beamer}
\usepackage{amsmath}
\usepackage[noend]{algorithmic}

\mode<presentation>{
\definecolor{cured}{rgb}{.8,0,.2}
\usecolortheme[named=cured]{structure}
\usetheme{split}
}

\usepackage{amsopn}

\newcommand{\GG}{\mathrm{GG}}
\newcommand{\YG}{\mathrm{YG}}
\newcommand{\RDG}{\mathrm{RDG}}

 
\usepackage{amsthm}

\newcommand{\centeripe}[1]{\begin{center}\Ipe{#1}\end{center}}
\newcommand{\comment}[1]{}

\newcommand{\centerpsfig}[1]{\centerline{\psfig{#1}}}

\newcommand{\seclabel}[1]{\label{sec:#1}}
\newcommand{\Secref}[1]{Section~\ref{sec:#1}}
\newcommand{\secref}[1]{\mbox{Section~\ref{sec:#1}}}

\newcommand{\alglabel}[1]{\label{alg:#1}}
\newcommand{\Algref}[1]{Algorithm~\ref{alg:#1}}
\newcommand{\algref}[1]{\mbox{Algorithm~\ref{alg:#1}}}

\newcommand{\applabel}[1]{\label{app:#1}}
\newcommand{\Appref}[1]{Appendix~\ref{app:#1}}
\newcommand{\appref}[1]{\mbox{Appendix~\ref{app:#1}}}

\newcommand{\tablabel}[1]{\label{tab:#1}}
\newcommand{\Tabref}[1]{Table~\ref{tab:#1}}
\newcommand{\tabref}[1]{Table~\ref{tab:#1}}

\newcommand{\figlabel}[1]{\label{fig:#1}}
\newcommand{\Figref}[1]{Figure~\ref{fig:#1}}
\newcommand{\figref}[1]{\mbox{Figure~\ref{fig:#1}}}

\newcommand{\eqlabel}[1]{\label{eq:#1}}
%\newcommand{\eqref}[1]{(\ref{eq:#1})}
\newcommand{\Eqref}[1]{Equation~(\ref{eq:#1})}

\newtheorem{thm}{Theorem}{\bfseries}{\itshape}
\newcommand{\thmlabel}[1]{\label{thm:#1}}
\newcommand{\thmref}[1]{Theorem~\ref{thm:#1}}

\newtheorem{lem}{Lemma}{\bfseries}{\itshape}
\newcommand{\lemlabel}[1]{\label{lem:#1}}
\newcommand{\lemref}[1]{Lemma~\ref{lem:#1}}

\newtheorem{cor}{Corollary}{\bfseries}{\itshape}
\newcommand{\corlabel}[1]{\label{cor:#1}}
\newcommand{\corref}[1]{Corollary~\ref{cor:#1}}

\newtheorem{obs}{Observation}{\bfseries}{\itshape}
\newcommand{\obslabel}[1]{\label{obs:#1}}
\newcommand{\obsref}[1]{Observation~\ref{obs:#1}}

\newtheorem{cond}{Condition}{\bfseries}{\itshape}

\newtheorem{clm}{Claim}{\bfseries}{\itshape}
\newcommand{\clmlabel}[1]{\label{clm:#1}}
\newcommand{\clmref}[1]{Claim~\ref{clm:#1}}


\newtheorem{dfn}{Definition}{\bfseries}{\rm}

\newtheorem{assumption}{Assumption}{\bfseries}{\rm}
\newenvironment{ass}{\begin{assumption}\rm}{\end{assumption}}
\newcommand{\asslabel}[1]{\label{ass:#1}}
\newcommand{\assref}[1]{Assumption~\ref{ass:#1}}

\newcommand{\proclabel}[1]{\label{alg:#1}}
\newcommand{\procref}[1]{Procedure~\ref{alg:#1}}

\newtheorem{rem}{Remark}
\newtheorem{op}{Open Problem}

\newcommand{\etal}{\emph{et al}}

%\newcommand{\keywords}[1]{\noindent\textbf{Keywords:} #1}
\newcommand{\voronoi}{Vorono\u\i}
\newcommand{\ceil}[1]{\left\lceil #1 \right\rceil}
\newcommand{\floor}[1]{\left\lfloor #1 \right\rfloor}
\newcommand{\R}{\mathbb{R}}
\newcommand{\N}{\mathbb{N}}
\newcommand{\Z}{\mathbb{Z}}
\newcommand{\Sp}{\mathbb{S}}
\newcommand{\E}{\mathrm{E}}

\usepackage{marvosym}

\newcommand{\notice}[1]
{
   {\Lightning}
   \marginpar{
      \begin{flushleft}\raggedright
        \hspace{-1.5mm}\Lightning{\small #1}
      \end{flushleft}
   }
}



\newcommand{\pin}{p_{\mathsf{in}}}
\newcommand{\pout}{p_{\mathsf{out}}}
\DeclareMathOperator{\cost}{cost}
\DeclareMathOperator{\depth}{depth}

\title{Moser's Proof of the Lov\'asz Local Lemma}
\author{Pat Morin}
\date{June 18, 2009 \\ SACT Seminar}

\begin{document}

\frame{\titlepage}

\frame
{
  \frametitle{STOC 2009}

  \begin{center}
    \includegraphics[height=1.3in]{moser} \\
    Robin A. Moser
  \end{center}
  \begin{itemize}
    \item PhD student at ETH Z\"urich
    \item Winner: STOC 2009 Best Student Paper Award
    \item Co-Winner: STOC 2009 Best Paper Award
  \end{itemize}
  R. A. Moser. A constructive proof of the Lov\'asz Local Lemma.
  \emph{Proceedings of STOC 2009}, pages~343--350. ACM Press, 2009. 
} 

%\section[Outline]{}
%\frame{\tableofcontents}

%\newcommand{\or}{\vee}
%\newcommand{\and}{\wedge}

\frame
{
  \frametitle{$k$-CNF Satisfiability}
  
  \begin{dfn}[$k$-CNF formula]
    A $k$-CNF formula is a conjunction of $m$ clauses, each of which is the
    disjunction of $k$ (possibly negated) distinct variables in 
    $x_1,\ldots,x_n$.
  \end{dfn}

  Example ($k=3$, $m=5$, $n=4$): 
  \begin{eqnarray*}
   F(x_1,\ldots,x_4) & = &
   (\bar x_1 \vee x_2 \vee x_3) 
   \wedge (x_1 \vee \bar x_2  \vee x_3) 
   \wedge (\bar x_1 \vee \bar x_2 \vee \bar x_3) \\
   & & {}\wedge (\bar x_1 \vee \bar x_3 \vee \bar x_4)
   \wedge (\bar x_1 \vee \bar x_2 \vee \bar x_4)
  \end{eqnarray*}
  \textbf{Quiz:} How long is the shortest unsatisfiable $k$-CNF formula?
}

\frame
{
  \frametitle{$k$-CNF Satisfiability}

  When is a $k$-CNF formula satisfiable? 

  \begin{itemize}
       \item NP-hard to decide, even for $k=3$ 
       \item Easy case: $F$ has less than $2^k$ clauses
       \item Easy case: Each variable appears in only one clause 
       \begin{itemize}
         \item Assign values $x_i\in \{0,1\}$ at random, so
         \[\Pr\{F(x_1,\ldots,x_n)=1\} = (1-(1/2)^k)^{m} > 0 \]
         \item There exists $x_1,\ldots,x_n$ such that $F(x_1,\ldots,x_n) = 1$.
       \end{itemize}
       \item What if each variable doesn't appear in more than $d$ clauses?
  \end{itemize} 
}


\frame
{
  \frametitle{The Lov\'asz Local Lemma}

  \begin{thm}[LLL --- Erd\"os and Lov\'asz 1975]
    Let $S=\{E_1,\ldots,E_m\}$ be a set of events such that
    \begin{enumerate}
      \item for all $i\in\{1,\ldots,m\}$, $\Pr\{E_i\} \le p$ 
      \item for all $i\in\{1,\ldots,m\}$, $E_i$ is dependent on at most $d$
events in $S\setminus \{E_i\}$, and 
      \item $4pd \le 1$.
    \end{enumerate}
    Then, $\Pr\{\bar E_1 \cap \bar E_2 \cap \cdots \cap \bar E_m\} > 0$.
  \end{thm}

  The LLL has applications in graph theory (esp. coloring), combinatorial geometry,
boolean satisfiability, network routing theory, scheduling theory, coding
theory, percolation theory, database theory, distributed computing, game
theory, wireless networks, Ramsey theory, network flow,\ldots

}

\frame
{
  \frametitle{Dependence}

  \begin{cond}[Limited Dependence]
    For each $i\in\{1,\ldots,m\}$, 
    there exists a \emph{dependent set} $S_i\subset S$ with
    \begin{enumerate}
      \item $|S_i| \le d$
      \item For any $S'\subseteq S\setminus S_i\setminus \{E_i\}$, $\Pr\{ E_i | \cap S_i'\} = \Pr\{E_i\}.$ 
    \end{enumerate}
   \end{cond}
}


\frame
{
  \frametitle{Using the LLL on $k$-CNF}

  \begin{thm}[$k$-CNF Sat]
    Let $F$ be a $k$-CNF formula where each variable appears in
    at most $d+1$ clauses, where $kd \le 2^{k-2}$.  Then $F$
    is satisfiable.
  \end{thm}
  \begin{proof}
    Assign $x_1,\ldots,x_n$ at random and let
    $E_i=$``the $i$th clause of $F$ is not satisfied''
    \begin{enumerate}
    \item $p=\Pr\{E_i\} = (1/2)^k$
    \item $E_i$ depends on at most $d'=kd$ other $E_j$ with $j\neq i$
    \item $4pd' \le 1$
  \end{enumerate}
  So, by the LLL $\Pr\{F(x_1,\ldots,x_n)=1\} > 0$ and therefore there
exists $x_1,\ldots,x_n$ such that $F(x_1,\ldots,x_n) = 1$.
  \end{proof}
}

\frame
{
  \frametitle{Constructive LLL}

  \begin{itemize}
    \item Theorem 1 proves the existence of a satisfying assignment, but
          doesn't explain how to find it
    \item For long formulas, a random assignment may have only a  very 
       small probability of working [$(1-(1/2)^{k-1}))^m$]
    \item We may have to try (exponentially) many random assignments before
       finding one that works
    \item Moser's new proof solves all this
    \begin{itemize}
      \item The proof is constructive
      \item The construction gives a randomized algorithm whose expected
            running time is polynomial in the number of clauses
    \end{itemize}
  \end{itemize}
}




\frame
{
  \frametitle{Moser's Algorithm}

  \noindent\textsc{Satisfy$(F=(C_1\wedge C_2\wedge\cdots\wedge C_m))$}
  \begin{algorithmic}
   \STATE{$x_1,\ldots,x_n \gets\mbox{ random element of }\{0,1\}^n$}
   \FOR{$i=1\ldots, m$}
     \IF{$C_i$ is false}
       \STATE{\textsc{Fix$(C_i)$}}
     \ENDIF
   \ENDFOR
  \end{algorithmic}
  \vspace{1ex}
  \noindent\textsc{Fix$(C_i)$}
  \begin{algorithmic}
   \STATE{Assign new random values to each of $C_i$'s $k$ variables}
   \FOR{each $C_j$ that shares a variable with $C_i$ (including $C_i$)}
     \IF{$C_j$ is false}
      \STATE{\textsc{Fix($C_j$)}}
     \ENDIF
   \ENDFOR
  \end{algorithmic}
}

\frame
{
  \frametitle{Communication Complexity}
 
  \begin{center} 
    \begin{tabular}{cc}
      \includegraphics[height=1in]{alice3} 
         & \includegraphics[height=1in]{bob-hope} \\
      Alice & Bob 
    \end{tabular}
  \begin{thm}[Shannon 1953]
    Let $C:\{0,1\}^r\rightarrow \{0,1\}^+$ be a variable-length code
    that encodes any $r$ bit string. If $B$ is a random $r$ bit string
    then $\E[|C(B)|] \ge r$.
  \end{thm}
  \end{center}
}

\frame
{
  \frametitle{Moser's Algorithm}
  \begin{itemize}
    \item We want to prove that Moser's Algorithm terminates, and in
polynomial time
    \item Let $B=b_1,b_2,b_3,\ldots$ be a random binary string
    \item When Moser's Algorithm needs a new random bit, it reads it from $B$
    \item Proof Idea:
    \begin{itemize}
      \item Moser's Algorithm encodes (a prefix of) $B$ very efficiently
      \item If it runs too long then the encoding violates Shannon's Theorem
        by encoding $r$ bits of $B$ using less than $r$ bits!
    \end{itemize}
  \end{itemize}
}

\frame
{
  \frametitle{Moser's Algorithm}

  \noindent\textsc{Satisfy$(F=(C_1\wedge C_2\wedge\cdots\wedge C_m))$}
  \begin{algorithmic}
   \STATE{$x_1,\ldots,x_n \gets b_1,\ldots,b_n$}
   \FOR{$i=1\ldots, m$}
     \IF{$C_i$ is false}
       \STATE{\textsc{Fix$(C_i)$}}
     \ENDIF
   \ENDFOR
  \end{algorithmic}
  \vspace{1ex}
  \noindent\textsc{Fix$(C_i)$}
  \begin{algorithmic}
   \STATE{Assign the next $k$ bits of $B$ to $C_i$'s $k$ variables}
   \FOR{each $C_j$ that shares a variable with $C_i$ (including $C_i$)}
     \IF{$C_j$ is false}
      \STATE{\textsc{Fix($C_j$)}}
     \ENDIF
   \ENDFOR
  \end{algorithmic}
}

\frame
{
  \frametitle{Tracing Moser's Algorithm}

  \noindent\textsc{SatisfyE$(F=(C_1\wedge C_2\wedge\cdots\wedge C_m))$}
  \begin{algorithmic}
   \STATE{$x_1,\ldots,x_n \gets b_1,\ldots,b_n$}
   \FOR{$i=1\ldots, m$}
     \IF{$C_i$ is false}
       \STATE{print $i$ and \textsc{FixE$(C_i)$} \COMMENT{* $\log m$ bits *}}
     \ENDIF
   \ENDFOR
  \end{algorithmic}
  \vspace{1ex}
  \noindent\textsc{FixE$(C_i)$}
  \begin{algorithmic}
   \STATE{Assign the next $k$ bits of $B$ to $C_i$'s $k$ variables}
   \STATE{$C_{i,1},\ldots,C_{i,\ell}\gets\{C_j:\mbox{$C_i$ and $C_j$ share a
variable}\}$}
   \FOR{$j=1,\ldots,\ell$}
     \IF{$C_{i,j}$ is false}
      \STATE{print $j$ and \textsc{FixE($C_{i,j}$)} \COMMENT{* $\log(kd)$ bits *}}
     \ENDIF
   \ENDFOR
   \STATE{print $\phi$ \COMMENT{* $O(1)$ bits *}}
  \end{algorithmic}
}


\frame
{
  \frametitle{Decoding Moser's Algorithm}

  \noindent$\textsc{SatisfyD}(F=(C_1\wedge C_2\wedge\cdots\wedge C_m))$
  \begin{algorithmic}
   \STATE{$x_1,\ldots,x_n \gets b_1,\ldots,b_n$}
   \WHILE{$i\gets\textsc{ReadInteger}$}
     \STATE{$\textsc{FixD}(C_i)$}
   \ENDWHILE
  \end{algorithmic}
  \vspace{1ex}
  \noindent$\textsc{FixD}(C_i)$
  \begin{algorithmic}
   \STATE{record values of variables in $C_i$ \COMMENT{* $k$ bits *}}
   \STATE{Assign the next $k$ bits of $B$ to $C_i$'s $k$ variables}
   \STATE{$C_{i,1},\ldots,C_{i,\ell}
        \gets\{C_j:\mbox{$C_i$ and $C_j$ share a variable}\}$}
   \WHILE{$j\gets \textsc{ReadInteger} \ne \phi$}
     \STATE{$\textsc{FixD}(C_{i,j})$}
   \ENDWHILE 
  \end{algorithmic}
}


\frame
{
  \frametitle{Tracing Moser's Algorithm}

  \begin{itemize}
    \item When \textsc{SatisfyD} completes, it has deduced all but 
          $n$ of the bits used by Moser's Algorithm
    \item Moser's Algorithm uses $n + sk$ bits of $B$, where $s$ is the
          number of calls to \textsc{Fix}.
    \item So, \textsc{SatisfyD} deduces $sk$ bits of $B$
  \end{itemize}
}

\frame
{
  \frametitle{Bounding Moser's Algorithm}

  \begin{itemize}
    \item \textsc{SatisfyE} prints $m\log m + s(\log(dk)+O(1))$ bits.
    \item By Shannon's Theorem, we must have
      \[ sk \le  m\log m + s(\log(dk)+O(1)) \]
      or
      \[ sk - s(\log(dk)+O(1)) \le m\log m \]
      \[ s(k - log(dk) - O(1)) \le m\log m \]
      \[ s \le \frac{m\log m}{k - \log(dk) - O(1)} \]
     provided that $k-\log(dk)-O(1) > 0$, i.e., $2^{k-O(1)} > dk$.
  \end{itemize}
}

\frame
{
  \frametitle{Moser's Theorem}

  \begin{thm}[Constructive $k$-CNF Sat]
    Let $F$ be a $k$-CNF formula where each variable appears in
    at most $d+1$ clauses, where $kd \le 2^{k-O(1)}$.  Then $F$
    is satisfiable and there exists a randomized algorithm that 
    finds a satisfying assignment in
    $O\left(m + \frac{m\log m}{k - \log(dk) - O(1)}\right)$ expected time.
  \end{thm}
}

\frame
{
  \frametitle{Summary}

  \begin{itemize}
    \item This gives an efficient algorithm for satisfying bounded
      occurence $k$-CNF formulas
    \item Moser and Tard\"os have a similar algorithm for the general LLL:
  \end{itemize}
  \begin{thm}[Lov\'asz and Tardos 2008]
    Let $\mathcal{P}$ be a finite set of mutually independent random
    variables in a probability space. Let $\mathcal{A}$ be a finite set of
    events determined by these variables. If there exists an assignment
    of reals $x : A \mapsto (0, 1)$ such that $\forall A \in \mathcal{A}
    : \Pr\{A\} \le x(A) \prod_{B\in \Gamma_{\mathcal{A}}(A)}
     (1 - x(B))$,
    then there exists a random algorithm that assigns values to the
    variables in $\mathcal{P}$ avoiding all the events in $\mathcal{A}$.
    This algorithm resamples an event $A \in \mathcal{A}$ at most
    an expected $x(A)/(1 - x(A))$ times before it finds such an
    assignment. 
%    Thus the expected total number of resampling steps is
%    at most $\sum_{A\in\mathcal{A}} x(A)/(1-x(A))$.
   \end{thm}
}

\end{document}

